\documentclass[11pt]{article}

\usepackage[margin=1in]{geometry}
\usepackage{amsmath, amssymb}
\usepackage{booktabs}
\usepackage[T1]{fontenc}

\title{Zeit Optimisation Model Note}
\author{Project Note}
\date{April 14, 2026}

\begin{document}
\maketitle

\section{Purpose}
Zeit plans one workweek by assigning tasks to a Monday--Friday calendar with fixed events,
30-minute granularity, hard deadlines, and finite working hours. The implementation uses a
CP-SAT model when OR-Tools is available and falls back to a deterministic greedy heuristic
when it is not.

\section{Planning Horizon}
Let
\[
D = \{1, \dots, 5\}
\]
be the workdays in the planning week and let each day contain
\[
H = 16
\]
half-hour slots corresponding to the window \(09{:}00\) to \(17{:}00\). The compressed weekly
slot set is
\[
S = \{0, \dots, 79\}.
\]
This compressed index is a modelling device: nights are omitted from the optimisation state,
even though the final UI renders the selected slots back onto a wall-clock calendar.

\section{Inputs}
For each task \(i \in T\), the model takes:
\begin{itemize}
  \item duration \(p_i\) measured in 30-minute slots,
  \item priority \(w_i \ge 0\),
  \item optional due date \(d_i\),
  \item hard-deadline flag \(h_i \in \{0,1\}\).
\end{itemize}

For each fixed event \(e \in E\), the model derives a busy interval on the compressed slot axis.
Those intervals are immutable and remove capacity from the feasible region.

\section{Decision Variables}
For each task \(i \in T\):
\begin{align*}
  x_i &\in \{0,1\} && \text{1 if task } i \text{ is scheduled,} \\
  s_i &\in \mathcal{S}_i && \text{start slot, drawn from the valid start domain,} \\
  f_i &= s_i + p_i && \text{finish slot.}
\end{align*}

The valid start set \(\mathcal{S}_i\) is precomputed so that a task remains inside one workday.
If \(h_i = 1\), starts that would force the task to finish after the hard deadline are removed
from \(\mathcal{S}_i\) before the solver runs.

For tasks with a soft due date, the CP-SAT model also introduces
\[
o_i \in \{0,1\},
\]
where \(o_i = 1\) means the scheduled task finishes on time.

\section{Constraints}

\subsection{Single Contiguous Block}
Each task can occupy at most one interval:
\[
[s_i, f_i)
\]
with fixed length \(p_i\). Task splitting is not allowed in the current version.

\subsection{No Overlap}
All scheduled task intervals and all fixed event intervals satisfy a global non-overlap condition:
\[
[s_i, f_i) \cap [s_j, f_j) = \emptyset \quad \forall i \ne j,
\]
and no scheduled task may intersect any fixed event interval. In OR-Tools this is expressed
with \texttt{AddNoOverlap} over optional task intervals plus constant event intervals.

\subsection{Hard Deadlines}
If task \(i\) has a hard deadline, then
\[
f_i \le d_i
\]
whenever \(x_i = 1\). If no feasible start remains after combining the workday boundary and the
busy-event exclusions, the task is declared unscheduled with reason \texttt{hard\_due\_conflict}.

\subsection{Capacity and Feasibility Filters}
Tasks longer than a full workday are removed before optimisation:
\[
p_i > H \implies i \text{ is unscheduled.}
\]
This is reported as \texttt{outside\_work\_window}.

\section{Objective}
The CP-SAT objective is a weighted lexicographic approximation implemented as a single linear
maximisation:
\[
\max \sum_{i \in T}
\left(
10000 \, w_i p_i x_i
+ 500 \, w_i o_i
+ (|S| - s_i) x_i
\right).
\]

The three components mean:
\begin{itemize}
  \item schedule high-priority, long-duration tasks first,
  \item prefer finishing due-dated work on time,
  \item break ties toward earlier starts in the week.
\end{itemize}

Because the coefficient on \(x_i\) dominates the others, the model first maximises useful
scheduled work, then uses due-date performance and earliness as secondary preferences.

\section{Greedy Fallback}
When OR-Tools is unavailable, tasks are sorted by:
\begin{enumerate}
  \item descending priority,
  \item hard due dates before soft due dates before undated work,
  \item earlier due dates,
  \item longer durations,
  \item title as a deterministic tiebreaker.
\end{enumerate}
The heuristic then places each task in the first feasible slot that avoids current busy intervals.
This does not guarantee optimality, but it preserves the same business rules and produces a stable
demo schedule.

\section{Current Modelling Limits}
\begin{itemize}
  \item Each task is a single contiguous block.
  \item Only one workweek is considered at a time.
  \item Context switching, travel, and preferred sequencing are not modelled yet.
  \item Priority is treated as a scalar weight rather than a richer utility curve.
\end{itemize}

\section{UI Interpretation}
The calendar UI should therefore display three distinct outcomes:
\begin{itemize}
  \item locked event slots that are fixed by the input calendar,
  \item scheduled task blocks placed by the optimiser and coloured by priority,
  \item unscheduled tasks with explicit infeasibility reasons.
\end{itemize}
This makes the optimisation result interpretable without exposing solver internals directly.

\end{document}
